\end{equation}
Thus \(A^{(ij)}(0)=A\) and \(|A^{(ij)}(\phi)_{ij}|=|a_{ij}|\) for every \(\phi\).
\end{definition}

\begin{definition}[Phase-restricted spectral set]
For an angular budget \(\delta\in[0,\pi]\), define
\begin{equation}
\Sigma^{\mathrm{phase}}_{ij}(A;\delta)
=
\bigcup_{|\phi|\leq\delta}
\spec\!\left(A^{(ij)}(\phi)\right).
\label{eq:phase-spectrum}
\end{equation}
\end{definition}

The ordinary \(\eps\)-pseudospectrum admits the perturbational definition
\[
\Lambda_\eps(A)
=
\bigcup_{\norm{E}_2\leq\eps}\spec(A+E).
\]
The phase-restricted set is a thin, structured slice of this larger region.

\begin{proposition}[Containment in the ordinary pseudospectrum]
For \(\delta\in[0,\pi]\),
\begin{equation}
\Sigma^{\mathrm{phase}}_{ij}(A;\delta)
\subseteq
\Lambda_\eps(A),
\qquad
\eps=2|a_{ij}|\sin(\delta/2).
\label{eq:containment}
\end{equation}
\end{proposition}

\begin{proof}
The perturbation in \cref{eq:entry-family} has rank one.  Its only nonzero entry is
\[
a_{ij}(e^{i\phi}-1),
\]
whose modulus is
\[
|a_{ij}|\,|e^{i\phi}-1|
=2|a_{ij}|\left|\sin\frac{\phi}{2}\right|.
\]
For a one-entry matrix, the spectral and Frobenius norms equal the modulus of that entry.  The maximum over \(|\phi|\leq\delta\) is the stated \(\eps\).  Every matrix in the phase family is therefore an admissible perturbation in \(\Lambda_\eps(A)\).
\end{proof}

The proposition gives the construction a precise place in pseudospectral theory.  It also clarifies what \cps{} does \emph{not} claim.  It does not approximate the entire pseudospectrum.  It trades completeness for interpretability.

\subsection{The velocity of an eigenvalue}

The local response of a simple eigenvalue has an exact formula.

\begin{proposition}[Eigenvalue phase velocity]
Let \(\lambda_k(\phi)\) be a simple eigenvalue of \(A^{(ij)}(\phi)\), with right eigenvector \(x_k\) and left eigenvector \(y_k\), normalized arbitrarily but with \(y_k^*x_k\neq0\).  Then
\begin{equation}
\frac{d\lambda_k}{d\phi}
=
 i\,a_{ij}(\phi)
\frac{\overline{y_{k,i}}x_{k,j}}{y_k^*x_k},
\qquad
a_{ij}(\phi)=|a_{ij}|e^{i(\arg a_{ij}+\phi)}.
\label{eq:eigenvalue-velocity}
\end{equation}
\end{proposition}

\begin{proof}
For a differentiable simple eigenpair, standard perturbation theory gives
\[
\lambda_k'(\phi)=\frac{y_k^*A'(\phi)x_k}{y_k^*x_k}.
\]
Here \(A'(\phi)=i a_{ij}(\phi)e_ie_j^*\), yielding \cref{eq:eigenvalue-velocity}.
\end{proof}

The numerator identifies the coupling's participation in the mode.  The denominator is the left-right eigenvector overlap.  When that overlap is small, the eigenvalue is ill-conditioned and the trajectory can move rapidly.  This is the local bridge from a visible loop to non-normal fragility.

\subsection{When a coupling cannot move the spectrum}

Not every off-diagonal entry matters spectrally.  The graph structure tells us when it can.

Associate to \(A\) a directed support graph with edge \(i\to j\) when \(a_{ij}\neq0\).

\begin{proposition}[Cycle-support criterion]
If the directed edge \(i\to j\) lies on no directed cycle of the support graph, then the characteristic polynomial of \(A\) is independent of \(a_{ij}\).  Consequently,
\[
\spec(A^{(ij)}(\phi))=\spec(A)
\quad\text{for all }\phi.
\]
\end{proposition}

\begin{proof}
Expand \(\det(\lambda I-A)\) as a sum over permutations.  A monomial containing the off-diagonal factor \(a_{ij}\) corresponds to a permutation with \(i\mapsto j\).  The cycle decomposition of that permutation must then contain a directed cycle beginning with \(i\to j\), followed by a directed path from \(j\) back to \(i\) through other selected off-diagonal factors.  If no such directed cycle exists in the support graph, no determinant monomial contains \(a_{ij}\).  The characteristic polynomial is therefore independent of that entry.
\end{proof}

